%&latex
\documentclass[a4paper]{article}

\frenchspacing

\usepackage[cp1250]{inputenc}
\usepackage[czech]{babel}

\usepackage{a4wide}
\usepackage{amsmath, amsthm, amssymb, amsfonts}
\usepackage[mathcal]{eucal}
\usepackage{graphicx}
\usepackage{url}
\usepackage{color}
\usepackage{wrapfig}
\usepackage{capt-of}
\usepackage{float}
\usepackage{ifthen}
\usepackage{subfig}
\usepackage[normalem]{ulem}


% sirka a vyska textu nastavena jako papir, vsechny okraje vynulovany a pridano 20pt na kazdou stranu
% horizontalni rozmery
\setlength{\textwidth}{\paperwidth}
\addtolength{\textwidth}{-40pt}
\addtolength{\hoffset}{-1in}
\addtolength{\hoffset}{20pt}
\setlength{\oddsidemargin}{0in}
\setlength{\marginparsep}{0in}
% vertikalni rozmery
\setlength{\textheight}{\paperheight}
\addtolength{\textheight}{-60pt}
\addtolength{\voffset}{-1in}
\addtolength{\voffset}{20pt}
\setlength{\topmargin}{0in}
\setlength{\headheight}{0in}
\setlength{\headsep}{0in}


%Obrazek na miste
%pouziti
%%\obrazeknahore{adresa}{popisek}{label}
\long\def\obrazeknahore#1#2#3 {

\begin{figure}[t]
    \centering
    \includegraphics[width=0.8\textwidth]{#1}
    
    \caption{#2}
    \label{#3}
    
\end{figure}

}


%==========================================
%PEKELNA MAKRA NA ZAROVNANI OBRAZKU DOPRAVA

\makeatletter


%tohle je makro, ktere mi dovoluje obtekani i u kratkych environmentu
%ABSOLUTNE nechapu, jak to funguje, ale funguje to
%viz http://tex.stackexchange.com/questions/26078/ 
\def\odrovnej{\@@par
\ifnum\@@parshape=\z@ \let\WF@pspars\@empty \fi % reset `parshape'
\global\advance\c@WF@wrappedlines-\prevgraf \prevgraf\z@
\ifnum\c@WF@wrappedlines<\tw@ \WF@finale \fi}

\makeatother



%---
%makro, co da obrazek doprava a ostatni text ho obteka
%(bez toho predchazejiciho makra to ale poradne nebeha)
%pouziti:
%\obrazekvpravo{adresa}{popisek}{label}{procento sirky}
\long\def\obrazekvpravo#1#2#3#4{

\setlength\intextsep{-20pt}

    \begin{wrapfigure}{r}{#4\textwidth}
      \begin{center}
          \vspace{-10pt}
          
        \includegraphics[width=#4\textwidth]{#1}
        \vspace{-10pt}
        
      \end{center}
      
      \caption{#2}
      \label{#3}
      
      
    \end{wrapfigure}

\setlength\intextsep{0pt}

    
}




%---
%makro pro pripady, kdy wrapfigure neco mrsi
%je to docela pekelne
%je nutne mu dat jak text vpravo, tak text vlevo
%a nevim, jestli bude 100% fungovat, ale doufam, ze jo

%pouziti:
%\obrazekvpravominipage{adresa}{popisek}{label}{procento sirky}{1 - procento sirky}{text vlevo}
\long\def\obrazekvpravominipage#1#2#3#4#5#6{

\noindent\begin{minipage}{#5\linewidth}
\vspace{0pt}
#6
\end{minipage}
\hspace{0.5cm}
\noindent\begin{minipage}{#4\linewidth}
\vspace{0pt}
\centering
\includegraphics[width=0.9\textwidth]{#1}
\captionof{figure}{#2}
\label{#3}
\end{minipage}

}

%KONEC PEKELNYCH MAKER
%=====================

\def\lebka{\includegraphics[width=3mm]{../common/lebka}}

% makra pro poznamku u vyrokove a predikatove logiky
\def\vl{ -- ve v�rokov� logice}
\def\pl{ -- v predik�tov� logice}


%Vacsina prostredi je dvojjazicne. V pripade, ze znenie napr pozorovania je pisane po slovensky, malo by byt po slovensky aj oznacenie.

\def\ifis#1#2{\ifthenelse{\equal{#1}{0}}{}{#2}}

%\newenvironment{e}[3]{\pagebreak[2]\noindent\textbf{\ifis{#2}{$\bigstar$}#1}\ifis{#3}{\emph{~(#3)}}\par\noindent\leftskip 10pt }{\odrovnej\par\bigskip}
\newenvironment{e}[3]{\pagebreak[2]\noindent\textbf{\ifis{#2}{\lebka~}#1}\ifis{#3}{\emph{~(#3)}}\ifis{#2}{~\lebka}\par\noindent\leftskip 10pt }{\odrovnej\par\bigskip}


%obecne prostredie, ktore ma vyuzitie pri specialnych odstavcoch ako (uloha, algoritmus...) aby nevzniklo dalsich x prostredi
\newenvironment{obecne}[1]{\pagebreak[2]\noindent\textbf{#1}\par\noindent\leftskip 10pt}{\odrovnej\par\bigskip}

\newenvironment{report}{\pagebreak[2]\noindent\textbf{Report}\em\par\noindent\leftskip 10pt}{\par\bigskip}

%\newenvironment{reportN}[1]{\pagebreak[2]\noindent\textbf{Report~}\emph{(#1)}\emph\par\noindent\leftskip 10pt}{\odrovnej\par\bigskip}
\newenvironment{reportN}[1]{\pagebreak[2]\noindent\textbf{Report~}\emph{(#1)}\em\par\noindent\leftskip 10pt}{\odrovnej\par\bigskip}

\newenvironment{penumerate}{
\begin{enumerate}
  \setlength{\itemsep}{1pt}
  \setlength{\parskip}{0pt}
  \setlength{\parsep}{0pt}
  %\setlength{\topsep}{200pt}
  \setlength{\partopsep}{200pt}
}{\end{enumerate}}

\def\pismenka{\numberedlistdepth=2} %pouzit, ked clovek chce opismenkovany zoznam...

\newenvironment{pitemize}{
\begin{itemize}
  \setlength{\itemsep}{1pt}
  \setlength{\parskip}{0pt}
  \setlength{\parsep}{0pt}
}{\end{itemize}}

\definecolor{gris}{gray}{0.95}
\newcommand{\ramcek}[2]{\begin{center}\fcolorbox{white}{gris}{\parbox{#1}{#2}}\end{center}\par}


\title{\LARGE U�ebn� texty k st�tn� bakal��sk� zkou�ce \\ Matematika \\ Posloupnosti a �ady funkc�}

\begin{document}

\maketitle

\newpage
\setcounter{section}{2}
\section{Posloupnosti a �ady funkc�}

\begin{e}{Po�adavky}{0}{0}
\begin{pitemize}
	\item Spojitost za p�edpokladu stejnom�rn� konvergence
	\item Mocninn� �ady
	\item Taylorovy �ady
	\item Fourierovy �ady
\end{pitemize}
\end{e}

\subsection{Stejnom�rn� konvergence}

\begin{obecne}{�mluva}
V dal��m budeme o $f$ a $f_n$ p�edpokl�dat, �e jsou definovan� na n�jak� nepr�zdn� mno�in� $M \subseteq \mathbb{R}$.
\end{obecne}

\begin{e}{Definice}{0}{t�i konvergence posloupnosti zobrazen�}
\begin{pitemize}
    \item �ekneme, �e posloupnost funkc� $(f_n)$ \emph{konverguje bodov� k funkci $f$ na mno�in� $M$} (zna��me $f_n \rightarrow f$ na $M$), jestli�e pro ka�d�  $x \in M$ plat� $\lim_{n \rightarrow \infty} f_n(x) = f(x)$. Explicitn� jestli�e
    $$\forall x \in M\, \forall \varepsilon > 0\; \exists n_0 \in \mathbb{N}:\; n \ge n_0 \Rightarrow |f_n(x)-f(x)| < \varepsilon.$$
    Tedy pro dan� $\varepsilon > 0$ m��e $n_0$ z�viset na bodu $x$.
    \item �ekneme, �e posloupnost funkc� $(f_n)$ \emph{konverguje stejnom�rn� k funkci $f$ na mno�in� $M$} (zna��me $f_n \rightrightarrows f$ na $M$), jestli�e
    $$\forall\varepsilon>0\ \exists n_0 \in \mathbb{N}:\; n\geq n_0, x\in M \Rightarrow |f_n(x) - f(x)| < \varepsilon.$$
    Zde naopak pro dan� $\varepsilon > 0$ mus� b�t $n_0$ univerz�ln� pro v�echny body $x$. �ekneme �e posloupnost funkc� je \emph{stejnom�rn� konvergentn�} na $M$, jestli�e konverguje k n�jak� funkci na $M$.
    \item �ekneme, �e posloupnost funkc� $(f_n)$ \emph{konverguje lok�ln� stejnom�rn� k funkci $f$ na mno�in� $M$} (zna��me $f_n \overset{loc}{\rightrightarrows} f$ na $M$), jestli�e pro ka�d� $x \in M$ existuje $\varepsilon > 0$ takov�, �e $f_n \rightrightarrows f$ na $M \cap (x-\varepsilon, x+\varepsilon)$. Tedy ka�d� $x \in M$ m� okol� (kter� m��e z�le�et na dan�m $x$), na n�m� $f_n \rightrightarrows f$.
\end{pitemize}
\end{e}

\begin{e}{D�sledek}{0}{0}
Z definic plyne s�la uveden�ch pojm�:
$$(f_n) \rightrightarrows f \text{ na } M \Rightarrow (f_n) \overset{loc}{\rightrightarrows} f \text{ na } M \Rightarrow (f_n) \to f \text{ na } M $$
\end{e}

\begin{e}{P��klad}{0}{0}
\begin{pitemize}
    \item $M = [0,1]$ a $f_n = x^n$, pak $(f_n) \to f$ spl�uj�c�
    $$f(x) =
    \begin{cases}
    0 &\text{ pro }x \in [0,1)\\
    1 &\text{ pro }x = 1.
    \end{cases}$$
    Ale nen� stejnom�rn� konvergentn�, jeliko� v bodech $1-1/n$ je posloupnost $(f_n)$ rostouc� a pro ka�d� $n \geq 2$ je $f_n(1-1/n) \geq 1/4$. Proto nen� ani lok�ln� stejnom�rn�, proto�e body $1-1/n$ konverguj� k $1$, a proto nem��e existovat ��dn� okol�, na n�m� by byla $(f_n)$ stejnom�rn� konvergentn�.\\
    Ov�em pro $M = [0, 1 - \delta]$ pro pevn� $\delta > 0$ m�me $0 \leq f_n(x) = x^n \leq (1-\delta)^n$. A jeliko� pro $n \to \infty$ jde $(1-\delta)^n \to 0$, dostaneme $|f_n(x) - f(x)| = |f_n(x)| \to 0$ nez�visle na $x \in [0, 1-\delta]$. Proto $(f_n) \rightrightarrows f \text{ na }M = [0, 1-\delta]$.\\
    Pro $M = [0,1)$ op�t kv�li bod�m $1-1/n$ m�me pouze lok�ln� stejnom�rnou konvergenci, jeliko� ka�d� bod $a \in [0,1)$ je obsa�en v okol� typu $[0, 1-\delta]$.
    \item $f_n(x) = \dfrac{nx}{1+n^2x^2}$ na $M = \mathbb{R}$ konverguje bodov� ke konstantn� funkci $f \equiv 0$. Body $x_n = 1/n$ jdou k nule, ale funk�n� hodnota z�st�v� $f_n(x_n) = 1/2$. Proto nejde ani o lok�ln� stejnom�rnou konvergenci. Kv�li tomu je $(f_n)$ na ka�d� $M$ neobsahuj�c� nulu lok�ln� stejnom�rn� konvergentn�. Kone�n� pro ka�dou $M$ neobsahuj�c� n�jak� okol� nuly z�sk�v�me stejnom�rnou konvergenci.
    \item $f_n(x) = \dfrac{\sin(nx)}{n} \rightrightarrows 0 \text{ na }M = \mathbb{R}$, jeliko� $|f_n(x)| \le 1/n$ pro v�echna $x \in \mathbb{R}$.
\end{pitemize}
\end{e}

\begin{e}{V�ta}{0}{spojitost limitn�ho zobrazen�}
Jsou-li $f_n$ spojit� a $f_n\rightrightarrows f$, pak je i $f$ spojit�.
\end{e}

\begin{e}{V�ta}{0}{Krit�rium stejnom�rn� konvergence}
Nech� $M$ je (nepr�zdn�) mno�ina, $f$ funkce definovan� na $M$ a $(f_n)$ posloupnost funkc� definovan�ch na $M$. Pak $f_n \rightrightarrows f$, pr�v� kdy�:
$$\lim_{n \to \infty} \sup_{x \in M}\left\lbrace\left\lvert f_n(x)-f(x)\right\rvert\right\rbrace = 0,$$
tj. existuje $n_0 \in \mathbb{N}$ takov�, �e pro $n \geq n_0$ je $\sup\limits_{x \in M}\left\lbrace\left\lvert f_n(x)-f(x)\right\rvert\right\rbrace$ definov�no (a kone�n�) a tato posloupnost m� limitu 0.
\end{e}

\begin{e}{P��klad}{0}{0}
\begin{pitemize}
    \item V metrick�m prostoru se supremovou metrikou je stejnom�rn� konvergence tot� jako bodov�.
    \item Metrick�m prostor omezen�ch funkc� na $M \subset \mathbb{R}$ se supremovou metrikou je �pln� a spojit� funkce v n�m tvo�� uzav�enou podmno�inu.
\end{pitemize}
\end{e}

\begin{e}{V�ta}{0}{Bolzanova-Cauchyova podm�nka pro stejnom�rnou konvergenci}
Posloupnost funkc� $(f_n)$ konverguje na mno�in� $M$ stejnom�rn� k n�jak� funkci $f$, pr�v� kdy�:
$$\forall \varepsilon > 0\, \exists n_0\in \mathbb{N}:\;m, n \ge n_0,\,x \in M \Rightarrow |f_n(x)-f_m(x)| < \varepsilon.$$
\end{e}

\begin{e}{Definice}{0}{monot�nn� bodov� konvergence}
�ekneme, �e bodov� konvergence $f_n \to f$ na $M$ je \emph{monot�nn�}, kdy� pro ka�d� bod $a \in M$ je posloupnost ��sel $(f_n(a))$ neklesaj�c� nebo kdy� pro ka�d� bod $a \in M$ je tato posloupnost nerostouc�.
\end{e}

\begin{e}{V�ta}{0}{0}
Plat� n�sleduj�c� tvrzen�:
\begin{penumerate}
    \item Kdy� $f_n \overset{loc}{\rightrightarrows}$ na $(a,b)$ pro $-\infty \leq a < b \leq \infty$, potom $f_n \rightrightarrows f$ na $[c,d]$ pro ka�d� kompaktn� podinterval $[c,d] \subset (a,b)$.
    \item (Diniho v�ta) Nech� $f_n \to f$ na kompaktn�m intervalu $I$, funkce $f_n$ i $f$ jsou spojit� a konvergence je monot�nn�. Pak $f_n \rightrightarrows f$ na $I$.
\end{penumerate}
\end{e}

\begin{e}{V�ta}{0}{Mooreova-Osgoodova o z�m�n� po�ad� limity $\lim_{n \to \infty}$ a $\lim_{x \to x_0}$}
Nech� $f$ i $f_n$ jsou funkce definovan� na n�jak�m prstencov�m okol� $P(x_0, \delta)$ bodu $x_0\in \mathbb{R}^*$, kter� m��e b�t i nevlastn�. D�le nech� $f_n \rightrightarrows f$ na $P(x_0, \delta)$ a pro ka�d� $n \in \mathbb{N}$ existuje vlastn� $\lim_{x \to x_0} f_n(x) = a_n$. Pak existuj� vlastn� limity $\lim_{n \to \infty} a_n$ a $\lim_{x \rightarrow x_0} f(x)$ a rovnaj� se:
$$\lim_{n \rightarrow \infty} \lim_{x \rightarrow x_0} f_n(x) = \lim_{x \rightarrow x_0} \lim_{n \rightarrow \infty} f_n(x).$$
\end{e}

\begin{e} %{Spojitost limitn� funkce}{D�sledek}{0}{0}
Nech� $I \subset \mathbb{R}$ je interval, $f_n$ spojit� funkce definovan� na $I$ a $f_n \overset{loc}{\rightrightarrows} f$ na $I$, pak $f$ je spojit� na $I$.
\end{e}

\begin{e}{V�ta}{0}{z�m�na po�ad� limity a integrov�n�}
Bu�te funkce $f_n$ a $f$ definovan� na (omezen�m) intervalu $[a, b]$, $f_n \in \mathcal{R}[a, b]$ a $f_n \rightrightarrows f$ na $[a, b]$. Pak i $f \in \mathcal{R}[a, b]$ a plat� $$ \int\limits_{a}^{b} f = \lim_{n \to \infty} \int\limits_{a}^{b} f_n.$$
\end{e}

\begin{e}{V�ta}{0}{z�m�na po�ad� limity a derivace}
Nech� $a, b \in \mathbb{R}, a<b$ a $(f_n)$ je posloupnost funkc� definovan�ch na intervalu $(a, b)$, kter� maj� v ka�d�m bod� $(a,b)$ vlastn� derivaci $f_n'$. Nech� d�le plat�:
\begin{penumerate}
	\item Existuje takov� $x_0 \in (a,b)$, �e posloupnost $\left(f_n(x_0)\right)$ je konvergentn�.
	\item Posloupnost $(f_n')$ je stejnom�rn� konvergentn� na $(a,b)$.
\end{penumerate}
Pak posloupnost $f_n \rightrightarrows f$ na $(a,b)$, funkce $f$ m� v ka�d�m bod� $x \in (a,b)$ vlastn� derivaci a plat� $f'(x) = \lim_{n \rightarrow \infty} f_n'(x)$.
\end{e}

\begin{obecne}{\color{red}TODO dal}
content...
\end{obecne}


\begin{e}{Definice}{0}{Bodov�/stejnom�rn� konvergence �ady funkc�}
�ekneme, �e �ada $\sum_{n=1}^{\infty} u_n$ \emph{konverguje bodov� na mno�in� $M$}, pokud posloupnost jejich ��ste�n�ch sou�t� je bodov� konvergentn� na $M$, tj, pro ka�d� $x \in M$ konverguje �ada $\sum_{n=1}^{\infty} u_n(x)$.

\emph{Sou�tem �ady} $\sum_{n=1}^{\infty} u_n$ nazveme funkci
$$S(x) = \sum_{n=1}^{\infty} u_n(x) = \lim_{n \rightarrow \infty} s_n(x), \, x \in M,$$
pokud �ada konverguje bodov� na $M$.

�ekneme, �e �ada $\sum_{n=1}^{\infty} u_n$ \emph{konverguje stejnom�rn� na mno�in� $M$}, pokud posloupnost jejich ��ste�n�ch sou�t� je stejnom�rn� konvergentn� na $M$.

Je-li nav�c $M \subset \mathbb{R}$, �ekneme, �e �ada $\sum_{n=1}^{\infty} u_n$ \emph{konverguje lok�ln� stejnom�rn� na mno�in� $M$}, pokud posloupnost jejich ��ste�n�ch sou�t� je lok�ln� stejnom�rn� konvergentn�.
\end{e}

\begin{e}{V�ta}{0}{Nutn� podm�nka stejnom�rn� konvergence �ady}
Nech� �ada $\sum_{n=1}^{\infty} u_n$ konverguje stejnom�rn� na mno�in� $M$. Pak $u_n \rightrightarrows 0$ na $M$.
\end{e}

\begin{e}{V�ta}{0}{Srovn�vac� krit�rium pro stejnom�rnou konvergenci}
Nech� M je (nepr�zdn�) mno�ina a $(u_n)_{n=1}^{\infty}$, $(v_n)_{n=1}^{\infty}$ dv� posloupnosti funkc� definovan�ch na M, pro kter� plat� $|u_n(x)| \le v_n(x)$ pro v�echna $x \in M$. Jestli�e �ada $\sum_{n=1}^{\infty} v_n$ konverguje stejnom�rn� na M, pak i �ada $\sum_{n=1}^{\infty} u_n$ konverguje stejnom�rn� na $M$.
\end{e}

\begin{e}{V�ta}{0}{Weierstrassovo krit�rium}
Nech� $M$ je (nepr�zdn�) mno�ina, $(u_n)_{n=1}^{\infty}$ posloupnost funkc� definovan�ch na $M$ a $\sum_{n=1}^{\infty} c_n$ konvergentn� �ada re�ln�ch ��sel. Pokud pro ka�d� $x \in M$ plat� $|u_n(x)| \le c_n$, pak �ada $\sum_{n=1}^{\infty} u_n$ konverguje stejnom�rn� na $M$.
\end{e}

\begin{e}{V�ta}{0}{Leibnizovo krit�rium pro stejnom�rnou konvergenci}
Nech� $M$ je (nepr�zdn�) mno�ina, $(u_n)_{n=1}^{\infty}$ posloupnost funkc� definovan�ch na $M$ spl�uj�c�ch \emph{ob�} podm�nky:

\begin{penumerate}
	\item Pro v�echna $x \in M$ a $n \in \mathbb{N}$ je $u_n(x) \ge u_{n+1}(x) \ge 0$
	\item $u_n \rightrightarrows 0$ na $M$
\end{penumerate}
Pak �ada $\sum_{n=1}^{\infty} (-1)^{n}u_n$ konverguje stejnom�rn� na $M$.
\end{e}

\begin{e}{V�ta}{0}{Dirichletovo a Abelovo krit�rium}
Nech� $M$ je (nepr�zdn�) mno�ina a $(u_n)_{n=1}^{\infty}$, $(v_n)_{n=1}^{\infty}$ dv� posloupnosti funkc� definovan�ch na $M$, p�i�em� pro ka�d� $x \in M$ a ka�d� $n \in \mathbb{N}$ plat� $v_n(x) \ge v_{n+1}(x) \ge 0$. Nech� nav�c plat� alespo� jedna z podm�nek:
\begin{enumerate}
	\item (Abelovo) �ada $\sum_{n=1}^{\infty} u_n$ konverguje stejnom�rn� na M, % a funkce $v_1$ je shora omezen� na $M$ -- to je malo, ne? (Tuetschek)
	pro ka�d� pevn� $x$ je posloupnost hodnot funkc� $(v_n(x))$ monot�nn� (klidn� pro ka�d� $x$ jinak)
	a existuje $K\in\mathbb{R}$ takov�, �e $\forall n\in\mathbb{N}\ \forall x\in M: |v_n(x)|<K$ 
	(tj. $(v_n)$ je stejnom�rn� omezen� na $M$). % tahle verze Abelova kriteria je podle Picka a Klazara
	\item (Dirichletovo) Existuje $K \in \mathbb{R}$ takov�, �e pro v�echna $x \in M$ a $n \in \mathbb{N}$ je $|u_1(x) + \dots + u_n(x)| \le K$
	(tj. posloupnost ��st. sou�t� $(\sum_{i=1}^n u_n(x))$ je \emph{stejnom�rn� omezen�} na $M$)
	a d�le $v_n \rightrightarrows 0$ na $M$ (konverguje stejnom�rn� k nulov� funkci).
\end{enumerate}

Pak �ada $\sum_{n=1}^{\infty} u_n\cdot v_n$ konverguje stejnom�rn� na $M$.
\end{e}

(\emph{Pozn. autora: D�le plat� i v�ty ekvivalentn� v�t�m o z�m�n� limit p�i posloupnostech\dots})

\newpage

T�to ot�zka je vypracovan� hlavne pod�a skr�pt prof. Kalendu, tak�e je mo�n� �e niektor� vety (napr. od prof. Pultra) bud� ma� in� znenie. Hlavne �as� o Fourierov�ch funkci�ch vyzer� by� predn�an� odli�ne (menej obecne)\dots ;-(
\begin{flushright}
\textit{andree}
\end{flushright}

\begin{obecne}{\color{red}TODO! cele zkontrolovat}
content...
\end{obecne}

\subsection{Mocninn� �ady}

\begin{e}{Definice}{0}{0}
Nech� $a \in \mathbb{R}$ a $\{c_n\}_{n=0}^{\infty}$ je posloupnost re�ln�ch ��sel. Nekone�nou �adu funkc� tvaru $\sum_{n=0}^{\infty} c_n(x-a)^n$ naz�v�me \emph{mocninnou �adou o st�edu $a$}.
\end{e}

\begin{e}{Definice}{0}{0}
Nech� $\sum_{n=0}^{\infty} c_n(x-a)^n$ je mocninn� �ada o st�edu $a$. Jej�m \emph{polom�rem konvergence} rozum�me ��slo
$$R = \sup\{r \in \left<0, +\infty\right); \sum_{n=0}^{\infty}|c_n|r^n \textit{konverguje}\}\textrm{,}$$
je-li uveden� mno�ina shora omezen�. Nen�-li shora omezen�, klademe $R = +\infty$.
\end{e}

\begin{e}{V�ta}{0}{0}
Nech� $\sum_{n=0}^{\infty} c_n(x-a)^n$ je mocninn� �ada o st�edu $a$ a $R$ jej� polom�r konvergence.
\begin{penumerate}
	\item Je-li $|x-a| < R$, pak �ada $\sum_{n=0}^{\infty} c_n(x-a)^n$ konverguje absolutn�;\\
		Je-li $|x-a| > R$, pak �ada $\sum_{n=0}^{\infty} c_n(x-a)^n$ diverguje.
	\item Je-li $r \in (0,R)$, pak �ada $\sum_{n=0}^{\infty} c_n(x-a)^n$ konverguje stejnom�rn� na mno�in� $\overline{B}(a,r) = \{x \in \mathbb{R}; |x-a|\le r\}=\left<a-r,a+r\right>$.
	\item �ada $\sum_{n=0}^{\infty} c_n(x-a)^n$ konverguje lok�ln� stejnom�rn� na mno�in� $B(a,R) = \{x \in \mathbb{R}; |x-a| < R\}$.
\end{penumerate}

Body 2. a 3. jsou vlastn� ekvivalentn�. Je-li $R=\infty$, pak �ada konverguje lok�ln� stejnom�rn� na cel�m $\mathbb{R}$. 
\end{e}

\begin{e}{Pozn�mka}{0}{0}
Mno�in� $B(a, R)$, kde $R$ je polom�r konvergence mocninn� �ady $\sum_{n=0}^{\infty} c_n(x-a)^n$, se ��k� \emph{kruh konvergence}.
\end{e}

\begin{e}{V�ta}{0}{V�po�et polom�ru konvergence}
Nech� $\sum_{n=0}^{\infty} c_n(x-a)^n$ je mocninn� �ada o st�edu $a$ a $R$ jej� polom�r konvergence.
\begin{penumerate}
	\item Jestli�e $L = \limsup_{n \rightarrow \infty} \sqrt[n]{|c_n|}$, pak
		$$R=\left\{ \begin{array}{ll} \frac{1}{L}, & L>0, \\ +\infty, & L=0 \end{array}\right.$$
	\item T�� vzore�ek plat�, je-li $L = \limsup_{n \rightarrow \infty} \left|\frac{c_{n+1}}{c_n}\right|$
\end{penumerate}
Prvn� bod plyne z Cauchyova odmocninov�ho krit�ria konvergence �ady, druh� z D'Alembertova pod�lov�ho krit�ria. Stejn� tvrzen� plat� i pro limity dan�ch v�raz� v p��pad�, �e existuj�.
\end{e}

\begin{e}{V�ta}{0}{...\uv{jen} pomocn� pro n�sleduj�c�}
Nech� $\sum_{n=0}^{\infty} c_n(x-a)^n$ je mocninn� �ada o st�edu $a$ a $R$ jej� polom�r konvergence. Pak i mocninn� �ady $\sum_{n=0}^{\infty} n.c_n(x-a)^{n-1}$ a $\sum_{n=0}^{\infty} \frac{c_n}{n+1}(x-a)^{n+1}$ maj� polom�r konvergence R.
\end{e}

\begin{e}{V�ta}{0}{Derivace a integrace mocninn� �ady}
Nech� $\sum_{n=0}^{\infty} c_n(x-a)^n$ je mocninn� �ada o st�edu $a$ a $R>0$ jej� polom�r konvergence. Definujme funkci $f(x) = \sum_{n=0}^{\infty} c_n(x-a)^n, x \in B(a,R)$. Pak plat�:
\begin{penumerate}
	\item Funkce $f$ je spojit� na $B(a, R)$.
	\item Funkce $f$ m� v ka�d�m bod� $x \in B(a,R)$ vlastn� derivaci a plat� $f'(x) = \sum_{n=0}^{\infty} n\cdot c_n(x-a)^{n-1}$.
	\item Funkce $F(x) = \sum_{n=0}^{\infty} \frac{c_n}{n+1}(x-a)^{n+1}$ je primitivn� funkc� k $f$ na $B(a, R)$.
\end{penumerate}
\end{e}


\subsection{Taylorovy �ady}

\begin{e}{Definice}{0}{0}
Nech� funkce $f$ m� v bod� $a$ derivace v�ech ��d�. Pak �adu $\sum_{n=0}^{\infty} \frac{f^{(n)}(a)}{n!}(x-a)^n$ naz�v�me \emph{Taylorovou �adou funkce $f$ o st�edu $a$ v bod� $x$}.
\end{e}

\begin{e}{Pozn�mka}{0}{0}
Nech� funkce $f$ m� v bod� $a$ derivace v�ech ��d� a $x \in \mathbb{R}$. Pak funkce $f$ je v bod� $x$ sou�tem sv� Taylorovy �ady o st�edu $a$, pr�v� kdy� $\lim_{n \rightarrow \infty} (f(x) - T_n^a(x))=0$.
\end{e}

\begin{e}{V�ta}{0}{0}
Nech� $x > a$ a funkce $f$ m� v ka�d�m bod� intervalu $\left<a,x\right>$ derivace v�ech ��d�. Jestli�e plat� podm�nka
\begin{pitemize}
	\item existuje $C \in \mathbb{R}$ takov�, �e pro ka�d� $t \in (a,x)$ a ka�d� $n \in \mathbb{N}$ je $|f^{(n)}(t)| \le C$,
\end{pitemize}
pak funkce $f$ je v bod� x sou�tem sv� Taylorovy �ady o st�edu $a$. Analogicky pro p��pad $x < a$.
\end{e}

\begin{e}{V�ta}{0}{0}
Nech� $\sum_{n=0}^{\infty} c_n(x-a)^n$ je mocninn� �ada o st�edu $a$ a $R>0$ jej� polom�r konvergence. Definujme funkci $f(x) = \sum_{n=0}^{\infty} c_n(x-a)^n, x \in B(a,R)$. Pak �ada $$\sum_{n=0}^{\infty} c_n(x-a)^n$$ je \emph{Taylorovou �adou funkce $f$} o st�edu $a$, tj. pro ka�d� $n \in \mathbb{N} \cup \{0\}$ plat� $c_n = \frac{f^{(n)}(a)}{n!}$.
\end{e}

\textbf{V�znam Taylorov�ch �ad}:
\begin{pitemize}
\item aproximace funkc� -- p��klady (Taylorovy �ady element�rn�ch funkc�):
	$$\forall x \in \mathbb{R}: \exp x=\sum_{k=0}^{\infty} \frac{1}{k!}x^k$$
	$$\forall x \in \mathbb{R}: \sin x=\sum_{k=0}^{\infty} \frac{(-1)^{k-1}}{(2k-1)!}x^{2k-1} \;\;\;\;\dots$$
\item zjednodu�en� d�kaz� -- p��klad (D�kaz binomick� v�ty): % Zdroj: Pultrova skripta
    \par\medskip
    Rozvineme funkci $f(x)=(1+x)^{\alpha}$ v okol� nuly. Indukc� lze ov��it, �e 
    $f^{(k)}(x)=\alpha(\alpha-1)\cdot\dots\cdot(\alpha-k+1)\cdot(1+x)^{\alpha-k}$. 
    Taylorova �ada funkce $f(x)=(1+x)^{\alpha}$ konverguje na $(-1,1)$ a je rovna hodnot� $(1+x)^{\alpha}$:
    $$(1+x)^{\alpha}=\sum_{k=0}^{\infty}\frac{\alpha(\alpha-1)\cdot\dots\cdot(\alpha-k+1)}{k!}x^k=\sum_{k=0}^{\infty}\binom{\alpha}{k}x^k$$
    a to d�v� binomickou v�tu.
\end{pitemize}

\subsection{Fourierovy �ady}

\subsubsection{Obecn� Fourierovy �ady}

\begin{e}{Definice}{0}{0}
Nech� $\{\varphi_n\}_{n=1}^{\infty}$ je posloupnost komplexn�ch funkc� na $\left<a,b\right>$, z nich� ��dn� nen� konstantn� nulov�. �ekneme, �e tato posloupnost tvo�� \emph{ortogon�ln� (kr�tce OG) syst�m na $\left<a,b\right>$}, jestli�e pro ka�d� dv� r�zn� $m, n\in \mathbb{N}$ plat�:
$$\int_a^b \varphi_m \overline{\varphi_n} = 0$$
Pokud nav�c
$$\int_a^b |\varphi_n|^2= 1$$
pro v�echna $n \in \mathbb{N}$, ��k�me, �e jde o \emph{ortonorm�ln� syst�m}.
\end{e}

\begin{e}{Pozn�mka}{0}{0}
P��klady OG syst�m�:
\begin{pitemize}
	\item Syst�m tvo�en� funkcemi $\exp \frac{2k\pi i x}{p}, k \in \mathbb{Z}$ je OG na intervalu $\left<a, a+p\right>$ pro ka�d� $a \in \mathbb{R}$
	\item Syst�m tvo�en� funkcemi $1, \cos \frac{2k\pi x}{p}, \sin \frac{2k\pi x}{p}, k \in \mathbb{N}$ je OG na intervalu \par$\left<a, a+p\right>$ pro ka�d� $a \in \mathbb{R}$
\end{pitemize}
\end{e}

\begin{e}{V�ta}{0}{0}
Nech� $\{\varphi_n\}_{n=1}^{\infty}$ je posloupnost komplexn�ch funkc� na $\left<a,b\right>$, $\{a_n\}_{n=0}^{\infty}$ je posloupnost komplexn�ch ��sel. Jestli�e
$$f(x) = \sum_{n=1}^{\infty} a_n \varphi_n(x), \, x \in \left<a,b\right>,$$
a uveden� �ada konverguje stejnom�rn� na $\left<a,b\right>$, pak pro ka�d� $n \in \mathbb{N}$ plat�
$$a_n = \frac{\int_a^b f\overline{\varphi_n}}{\int_a^b |\varphi_n|^2}.$$
\end{e}

\begin{e}{Definice}{0}{po ��stech spojit� funkce}
�ekneme, �e funkce $f$ je \emph{po ��stech spojit� na $\left<a,b\right>$}, jestli�e existuje $D=\{x_i\}_{j=0}^{N}$ d�len� intervalu $\left<a,b\right>$ takov�, �e pro ka�d� $j \in \{1,\dots,N\}$ je funkce $f$ spojit� na intervalu $(x_{j-1}, x_j)$ a v krajn�ch bodech tohoto intervalu m� vlastn� jednostrann� limity.
\end{e}

\begin{e}{Definice}{0}{0}
Nech� $\{\varphi_n\}_{n=1}^{\infty}$ je OG syst�m na $\left<a,b\right>$ a funkce $f$ je po ��stech spojit� na $\left<a,b\right>$. Pro $n \in \mathbb{N}$ polo�me
$$a_n = \frac{\int_a^b f\overline{\varphi_n}}{\int_a^b |\varphi_n|^2}.$$
Tato ��sla naz�v�me \emph{Fourierov�mi koeficienty funkce $f$ vzhledem k OG syst�mu $\{\varphi_n\}_{n=1}^{\infty}$ na $\left<a,b\right>$} a �adu
$$\sum_{n=1}^{\infty} a_n \varphi_n$$
naz�v�me \emph{Fourierovou �adou $f$ vzhledem k OG syst�mu $\{\varphi_n\}_{n=1}^{\infty}$ na $\left<a,b\right>$}.
\end{e}

\subsubsection{Trigonometrik� Fourierovy �ady}

\begin{e}{Definice}{0}{po ��stech spojit� periodick� funkce}
Bu� funkce $f$ periodick� s periodou $p > 0$. �ekneme, �e je \emph{po ��stech spojit�}, je-li po ��stech spojit� na intervalu $\left<0,p\right>$.
\end{e}

\begin{e}{Pozn�mka}{0}{0}
Nech� $f$ je $p$-periodick� funkce a $a,b�\in \mathbb{R}$.
\begin{penumerate}
	\item Pak $f$ je po��stech spojit� na $\left<a, a+p\right>$, pr�v� kdy� je po ��stech spojit� na $\left<b, b+p\right>$.
	\item $\int_a^{a+p}f = \int_b^{b+p}f$, pokud alespo� jeden z t�chto integr�l� existuje.
\end{penumerate}
\end{e}

\begin{e}{Definice}{0}{0}
Nech� funkce $f$ je $p$-periodick� po ��stech spojit� funkce. Jej�mi \emph{trigonometrick�mi Fourierov�mi koeficienty} rozum�me ��sla
$$a_n = \frac{2}{p} \int_0^p f(x)\cos \frac{2\pi nx}{p}dx, \, n \in \mathbb{N} \cup \{0\}$$
$$b_n = \frac{2}{p} \int_0^p f(x)\sin \frac{2\pi nx}{p}dx, \, n \in \mathbb{N}$$
\end{e}

\begin{e}{Definice}{0}{0}
\emph{Trigonometrickou Fourierovou �adou} funkce $f$ pak rozum�me �adu
$$\frac{a_0}{2} + \sum_{n=1}^{\infty} \left( a_n \cos \frac{2\pi nx}{p} + b_n \sin \frac{2\pi nx}{p}\right)$$
\end{e}

\begin{e}{Pozn�mka}{0}{Besselova nerovnost}
Besselova nerovnost pro trigonometrick� Fourierovy �ady m� tvar
$$\frac{|a_0|^2}{4}p + \sum_{n=1}^{\infty} (|a_n|^2 + |b_n|^2)\frac{p}{2} \le \int_0^p |f|^2.$$
Podobn� nerovnost plat� i pro obecn� Fourierovy �ady.
\par\medskip\noindent
(Riemann-Lebesgue) d�sledkem t�to nerovnosti je fakt, �e $\lim a_n = \lim b_n = 0$.
\end{e}

\begin{e}{V�ta}{0}{Persevalova rovnost}
Pro trigonometrick� Fourierovy �ady plat� v Besselov� nerovnosti rovnost. Pro funkce s periodou $2\pi$ potom plat�:
$$\frac{1}{\pi} \int_{-\pi}^{\pi} |f|^2= \frac{|a_0|^2}{2}+\sum_{n=1}^{\infty} (|a_n|^2 + |b_n|^2) \,\,\,\textit{(jedna z variant z�pisu)}$$
\end{e}

\begin{e}{Pozn�mka}{0}{0}
Nech� $f$ je $p$-periodick� po ��stech spojit� funkce takov�, �e v�echny jej� trigonometrick� Fourierovy koeficienty jsou nulov�. Pak $f(x)=0$ pro v�echna $x \in \left<0,p\right>$ s v�jimkou kone�n� mnoha bod�.
\end{e}

\begin{e}{V�ta}{0}{Symetrie funkce a Trigonometrick� Fourierovy koeficienty}
Nech� $f$ je $p$-periodick� po ��stech spojit� funkce, $a_n, n \in \mathbb{N} \cup \{0\}$ a $b_n, n \in \mathbb{N}$, jej� trigonometrick� Fourierovy koeficienty. Pak plat�
\begin{penumerate}
	\item Pro v�echna $n \in \mathbb{N} \cup \{0\}$ je $a_n=0$, pr�v� kdy� $f(-x) = -f(x)$ pro v�echna $x \in \left<0,p\right>$ s v�jimkou kone�n� mnoha bod�.
	\item Pro v�echna $n \in \mathbb{N}$ je $b_n=0$, pr�v� kdy� $f(-x) = f(x)$ pro v�echna $x \in \left<0,p\right>$ s v�jimkou kone�n� mnoha bod�.
\end{penumerate}
\end{e}

\begin{e}{Definice}{0}{0}
Nech� $f$ je $p$-periodick� po ��stech spojit� funkce. �ekneme, �e $f$ je \emph{po ��stech hladk�}, jestli�e $f'$ je po ��stech spojit�.
\end{e}

\begin{e}{V�ta}{0}{O konvergenci Fourierov�ch �ad}
Nech� $f$ je po ��stech hladk� p-periodick� funkce. Pak plat�:
\begin{penumerate}
	\item Trigonometrick� Fourierova �ada funkce $f$ konverguje bodov� na $\mathbb{R}$ a jej� sou�et v bod� $x \in \mathbb{R}$ je $\frac{1}{2} \left( \lim_{t\rightarrow x-} f(t) + \lim_{t\rightarrow x+} f(t)\right)$
	\item Je-li $f$ nav�c spojit� na intervalu $(a,b)$, pak jej� trigonometrick� Fourierova �ada konverguje lok�ln� stejnom�rn� na $(a,b)$ a jej� sou�et je $f(x)$ pro ka�d� $x \in (a,b)$.
	\item Je-li nav�c spojit� na $\mathbb{R}$, pak jej� trigonometrick� Fourierova �ada konverguje stejnom�rn� na $\mathbb{R}$ a jej� sou�et je $f(x)$ pro ka�d� $x \in \mathbb{R}$.
\end{penumerate}
\end{e}

% derivacia FR?: podle me netreba, ani Pick ani Klazar ani Pultr to nemaji ve skriptech / neprednaseli -- Tuetschek


\end{document}
